\chapter{Conclusions and Future Work}
\label{ch:conclusions}

% Draft written with AI assistance (12 Nov 2026) for the author to check and revise.
% Every number is the one reported in Chapters 5 and 6.

\section{Conclusions}
\label{sec:conclusions}

This thesis examined how well the emotions that listeners perceive in film music predict the genre of the film, and how an interpretable emotion step compares with classifying genre
directly from audio. The emotion pipeline predicts eight emotion ratings from a VGGish
embedding with a random forest and then predicts genre from these ratings with one
logistic regression per genre, so that a film can receive several genres (multi-label
classification). It was evaluated on the 40 films of the five-genre Eerola
subset and transferred to the 110 films of the Blockbuster dataset. Section~\ref{sec:disc_rq}
answers the four research questions in detail; in brief, the results are as follows.

The emotions rank films above chance (macro ROC-AUC 0.655), but at the default threshold their macro-$F_1$ is close to that of a trivial classifier. The emotion pipeline reaches a film-level macro-$F_1$ of 0.460 and the human ratings 0.491, against 0.309 for a random guess and 0.434 for assigning every genre to every film. Most genres have an emotional signature that is consistent with the conventions of
film scoring. Fear is the most discriminating emotion and valence a second axis, and one to
three emotions keep most of the genre information. Within one dataset, the emotion pipeline
is ahead of every audio representation used directly at the default threshold ($+0.093$
over VGGish, $p=0.052$), but significantly only against the weaker ones, and with
thresholds tuned per genre the pipelines converge. Trained on Eerola alone, it transfers to
the Blockbuster films better than direct VGGish at the default threshold (0.508 against
0.403, $p=0.013$), but most of this lead depends on the threshold.

\paragraph{Overall conclusion.} The results do not indicate that the emotion step makes
genre prediction from film music more accurate in general. Its lead appeared on the small dataset and under transfer, at the default decision threshold, and much of it depends on that threshold; whether the amount of training data is the cause was not tested. Its main value lies in interpretability: its predictions rest on eight emotions that listeners can name, and some errors can be traced to a stage of the pipeline. The weakness on Horror, for
example, can be traced mainly to Stage~1, which shifts the extreme fear ratings of Horror clips towards the
average (rated 4.35, predicted 3.73; Appendix~\ref{sec:app_horror}). Such an explanation can
be checked against the practice of film scoring, which is not possible in the same way for a
direct audio model.

\section{Future Work}
\label{sec:futurework}

Several experiments could build on this thesis. Some would test the present results on more
data, some would transfer the emotion step to new kinds of data, and some would apply it in
other domains.

\paragraph{A larger replication with fixed decisions.} The most direct next step is a
replication on a larger collection of film-music clips with emotion ratings and genre
labels. The unit of evaluation, the genre set, the threshold policy and the main comparisons
should be fixed before the data are analysed. Such a study could test two leads that this
thesis could only observe: the advantage of the emotion pipeline over the strongest direct
representations, and the lead of CLAP with ridge regression, which was found after the
results were known ($+0.032$, $p=0.112$; Appendix~\ref{sec:app_grid}). The same collection
could include a genre judgement for every single clip, which would separate errors of the
model from the fact that IMDb labels describe whole films.

\paragraph{Transfer with audio and ratings on both sides.} The Blockbuster dataset has
neither audio nor emotion ratings, so the transfer in this thesis was limited to VGGish and
could not check the predicted emotions. A second dataset with audio and its own ratings would
allow Stage~1 to be tested on new films and any representation to be transferred. It would
also allow a few rated clips of the new dataset to be used to adapt Stage~1 before genre is
predicted. Comparing this adaptation with zero-shot transfer would show how much labelled
data a new collection needs.

\paragraph{New sources of film music.} The emotion step could be transferred to other kinds
of audiovisual music. Trailers are a natural target, because their music is edited to
advertise the genre, and trailer audio has already been used for genre classification
\citep{mangolin2020multimodal}. Films from other film industries and periods would show
whether the emotional conventions of genres found here hold beyond the films of this dataset.
Television series and video games use similar conventions and could be tested in the same
way.

\paragraph{A stronger emotion model.} Stage~1 reaches about 62\,\% of the attainable $R^2$
and underestimates the strongest emotions, which appears to cause much of the weakness on
Horror. Audio models fine-tuned on emotion ratings, instead of frozen embeddings, could be
compared with the random forest. Each such model should be evaluated on genre as well,
because a more accurate emotion model did not reliably yield a better genre result in this
thesis (Section~\ref{subsec:eval_grid}).

\paragraph{Calibration and thresholds as part of the design.} Much of the advantage of the
emotion pipeline depends on the default threshold of 0.5. A future study could compare
calibrated probabilities with thresholds chosen per genre \citep{lipton2014optimal} and
report a threshold-free measure such as ROC-AUC beside macro-$F_1$ from the start.

\paragraph{Emotion in multimodal models.} Audio alone is a weak source of genre information
compared with images and text \citep{mangolin2020multimodal}. The emotion step could serve
as an interpretable audio component beside trailer frames or the plot synopsis
\citep{behrouzi2023multimodal}. Such a model would show whether the emotions still add
information when stronger sources are available, and it would allow a genre prediction to
be explained partly in terms of the music.

\paragraph{Applications in other domains.} An emotion step between audio and a category
could be useful wherever music has to fit a purpose. Examples are the search for music in
production libraries for a scene of a given mood or genre, music recommendation that can
state why a piece was chosen, and the tagging of music for advertising or games. In each
case the emotions would make the decision readable for the user, which this thesis suggests
is their main strength.
